\documentclass[12pt]{jlreq}
\usepackage{amsmath, amssymb}

\newcommand{\C}{\mathbb{C}}
\newcommand{\R}{\mathbb{R}}
\newcommand{\Z}{\mathbb{Z}}
\newcommand{\N}{\mathbb{N}}
\newcommand{\F}{\mathbb{F}}
\newcommand{\Prob}{\mathbb{P}}
\newcommand{\Exp}{\mathbb{E}}
\newcommand{\dd}{\,d}
\newcommand{\Ker}{\operatorname{Ker}}
\newcommand{\Impart}{\operatorname{Im}}
\newcommand{\Hom}{\operatorname{Hom}}

\begin{document}

\[
  A:=\C[x],\qquad
  B:=\C[x,y,z]/(y(xy-z^2),\,z(xy-z^2))
\]
を可換環とする。環準同型 \(f:A\to B\) を，単射準同型
\[
  \iota:\C[x]\longrightarrow \C[x,y,z],\qquad f(x)\longmapsto f(x)
\]
と，\(\C[x,y,z]\) のイデアル \((y(xy-z^2),\,z(xy-z^2))\) によって引き起こされる商準同型
\[
  \pi:\C[x,y,z]\longrightarrow
  \C[x,y,z]/(y(xy-z^2),\,z(xy-z^2))
\]
の合成で，\(f=\pi\circ\iota\) と定める。

\begin{enumerate}
  \item \(A\) の極大イデアルをすべて求めよ。説明も書くこと。
  \item \(B\) を \(f\) によって \(A\)-加群と見なすとき，\(A\) の各極大イデアル \(\mathfrak m\) に対して，
  \[
    C^{\mathfrak m}:=B\otimes_A A/\mathfrak m
  \]
  とおく。\(C^{\mathfrak m}\) を \(\C\) 上の \(2\) 変数多項式環の商環の形で書け。
  \item \(C^{\mathfrak m}\) の各極大イデアル \(\mathfrak n\) による局所化を \(C^{\mathfrak m}_{\mathfrak n}\) とおく。\(C^{\mathfrak m}_{\mathfrak n}\) の極大イデアル \(\overline{\mathfrak n}\) に対して
  \[
    \dim_{\C}\overline{\mathfrak n}^2/\overline{\mathfrak n}^3
  \]
  を計算せよ。
\end{enumerate}

\end{document}